% should this be in this section or the intro

Modeling the interaction between a laser beam and droplet combines the approaches of multiphase flow and magnetohydrodynamics (MHD). The former is simply due to the existence of the two immiscible liquid and gas phases,
and the latter is especially necessary if the laser intensity is strong enough to ionize. In this paper, the term gas phase includes both neutral and charged particles, as both are governed by the ideal gas equation of
state.

In order to accurately capture the behavior of the liquid-gas-plasma, a four-fluid model is used to describe the constituent fluids -- namely, liquid, gas, electrons and ions -- that interact with one another. The charged
particles also exhibit electromagnetic effects. The governing equations are therefore derived from the three conservation laws and Maxwell's equations.

Chemically speaking, the laser can depart enough energies to not only remove bound electrons from their orbits but also break molecular bonds and alter the chemical composition of the fluid. In fact, free radicals and
light molecular ions are frequently encountered in plasmas if the laser is not strong enough to completely disintegrate molecules into monoatomic ions. However, for simplicity, all ions are assumed to have resulted
from neutral molecules each losing a single electron, and thus they have a charge of +1 and the same molecular weight as the parent neutral particle.